\section{Introduction}
\label{sec:intro}

Post-training for long-horizon agentic capabilities has a tension between efficiency and accuracy. By long-horizon, we mean tasks that require many turns of interaction with an environment, such as conversational tool use, code editing, terminal interaction, and web search. Supervised fine-tuning (SFT) is efficient, but it often degrades out-of-domain (OOD) accuracy. End-to-end reinforcement learning (E2E RL) preserves on-policy behavior and typically retains OOD performance better, but it is expensive because every update requires repeated full multi-turn environment interaction and can therefore be slow in wall-clock time. This raises a natural question: can we recover the accuracy and OOD retention of E2E RL without paying the full-trajectory rollout cost?

A natural first attempt is to reuse expert trajectories for RL. Given an SFT expert trajectory, we can sample an assistant turn, condition on the expert prefix, and roll out only the next assistant action rather than the full trajectory. This idea is appealing because it converts existing SFT data into RL data and reduces online interaction from full trajectories to one-turn. However, the naive version of this idea does not work well in practice. Randomly chosen turns often provide little or no learning signal, and just exact expert matching is too strict in generative action spaces, where many functionally appropriate actions differ from the single demonstrated completion.

Our core claim is that both issues matter. First, many intermediate turns are non-informative under group-normalized RL, such as GRPO: if sampled actions at a turn all succeed or all fail, the normalized advantage is zero and no update is produced. Second, strict matching misses credit for acceptable alternatives: a tool call, shell command, or search action may be functionally appropriate for that point in the trajectory without matching the expert string exactly. Such a method must therefore answer two questions at once: which turns are worth training on, and how should success be scored at those turns?

PivotRL makes these two requirements explicit. Demonstrations expose many intermediate decision states, but behavior cloning underuses them because it trains on only one demonstrated completion at each state. PivotRL performs local on-policy exploration around these states. It filters to what we call pivots---assistant turns whose sampled outcomes offer mixed outcomes ---and it optimizes a functional reward that credits any action judged appropriate by a domain verifier. Training then requires only one-turn or partial rollouts from those pivots, rather than full online trajectories.

We provide a lightweight analysis for both components. A simple proposition explains why mixed-outcome turns are the only ones that produce nonzero group-normalized updates, which motivates pivot selection. A KL-projection theorem explains why the functional reward induces a conservative update relative to the reference policy: it increases total probability on verifier-approved actions while preserving the reference ordering within and outside that set. We present this analysis as support for the method, not as a complete theory of forgetting.

We evaluate PivotRL in three ways. First, we compare it to SFT under identical training data---the same prompts and expert trajectories---and show larger in-domain gains with substantially smaller OOD regression. Across four training domains, PivotRL achieves an average in-domain gain of +14.11 points relative to the base model, compared with +9.94 for SFT, while average OOD change is +0.20 for PivotRL versus $-9.48$ for SFT. Second, on SWE-Bench, where E2E RL is a natural baseline, we show that PivotRL attains competitive accuracy with much lower rollout burden. Third, we ablate the two components of the method and show that both informative pivot selection and the functional reward are necessary for the strongest results.

